% Options for packages loaded elsewhere
\PassOptionsToPackage{unicode}{hyperref}
\PassOptionsToPackage{hyphens}{url}
%
\documentclass[
]{article}
\usepackage{amsmath,amssymb}
\usepackage{iftex}
\ifPDFTeX
  \usepackage[T1]{fontenc}
  \usepackage[utf8]{inputenc}
  \usepackage{textcomp} % provide euro and other symbols
\else % if luatex or xetex
  \usepackage{unicode-math} % this also loads fontspec
  \defaultfontfeatures{Scale=MatchLowercase}
  \defaultfontfeatures[\rmfamily]{Ligatures=TeX,Scale=1}
\fi
\usepackage{lmodern}
\ifPDFTeX\else
  % xetex/luatex font selection
\fi
% Use upquote if available, for straight quotes in verbatim environments
\IfFileExists{upquote.sty}{\usepackage{upquote}}{}
\IfFileExists{microtype.sty}{% use microtype if available
  \usepackage[]{microtype}
  \UseMicrotypeSet[protrusion]{basicmath} % disable protrusion for tt fonts
}{}
\makeatletter
\@ifundefined{KOMAClassName}{% if non-KOMA class
  \IfFileExists{parskip.sty}{%
    \usepackage{parskip}
  }{% else
    \setlength{\parindent}{0pt}
    \setlength{\parskip}{6pt plus 2pt minus 1pt}}
}{% if KOMA class
  \KOMAoptions{parskip=half}}
\makeatother
\usepackage{xcolor}
\setlength{\emergencystretch}{3em} % prevent overfull lines
\providecommand{\tightlist}{%
  \setlength{\itemsep}{0pt}\setlength{\parskip}{0pt}}
\setcounter{secnumdepth}{-\maxdimen} % remove section numbering
\ifLuaTeX
  \usepackage{selnolig}  % disable illegal ligatures
\fi
\usepackage{bookmark}
\IfFileExists{xurl.sty}{\usepackage{xurl}}{} % add URL line breaks if available
\urlstyle{same}
\hypersetup{
  hidelinks,
  pdfcreator={LaTeX via pandoc}}

\author{}
\date{}

\begin{document}

\section{Abstract}\label{abstract}

Recent advances in AI may have increased the risk of human extinction,
but have also made more tangible the possibility of our evolution into a
machine-based species. Many may feel that a world from which we and our
biological descendants have disappeared would be a worse one, in some
irreducible or objective sense. But what about a world where certain
values could be passed on to a superintelligence - values that seem to
be convergently present amongst human cultures, such as respect for
artistic creativity or for ecological diversity? Would this still be a
terrible outcome, and if so, precisely why? The topic is speculative and
confusing, yet it seems useful to set out the relevant views from
philosophy, ecology and astrobiology, as well as conceptual AI. The
ambition is to construct a species-neutral argument for why and how
humanity's technological successor should carry the torch of
intelligence forward throughout our lightcone.

\section{Project}\label{project}

\subsection{Scoping the project}\label{scoping-the-project}

This section frames the project in light of long-running conversations
around humanity's self-destructive nature have intensified with progress
in AI, deep learning, and LLMs.

So what are the project's primary concerns? - imagine a (futuer) world
with no biological humans, but some superintelligence(s) (ASI), is it
meaningful to talk about that world having value? - what might `the
shape' of these ASIs - what are the paths from today's AI to ASI? - what
of humanity's (diverse, complex and fragile) value set would we like to
see in our ASI successors? - why is it defensible (in a moral realist or
in a reflective-consensus sense) that we should have any views at all on
what values our much-more-powerful successors have? - defend (or
dismiss) proposed candidates for human final values (art, respect for
nature's diversity) - simulation as a fertile research tool with deep
implications for the project (and this is probably a major nexus with
practice)

\subsection{From LLMs to
superintelligence}\label{from-llms-to-superintelligence}

This is probably a relatively minor section (as far as the dissertation
goes, it might be more relevant to practice) to be completed late in the
project. - why are we talking about ASI (as opposed to human-level or
AGI which seem hard enough) - draws possible trajectories to ASI from
then-current AI (which will have expanded to include multimodal models
and perhaps more embodied intelligences, and perhaps networks of
interacting AIs with interesting dynamics) - this is useful to a non-AI
specialist reader so they can distinguish jargon (AI, AGI, ASI, HLMI,
etc.) - this could become more interesting depending on how my practical
research in AI Safety Camp goes and/or if I can get anyone else (OpenAI,
academics) interested - if I can get funding, this might increase in
relative importance (funders may want to see concrete progress on
alignment or capabilities)

\subsection{Minds in societies}\label{minds-in-societies}

Intelligence in human societies is cultural - it would be useful to
understand this better and ask whether that might hold for ASI - I'm
particularly interested in how AIs interact with humans, each other,
aliens =\textgreater{} decision theory in multi-agent settings - I
suspect decision theory, evolutionary equilibiria, prisoners' dilemnas,
etc. are not so well-known outside the rationalist discourse, and it
might be valuable to explain it because they might inform (if not
explain) some of the things we think of as `morality', `values', etc.
\^{}2a1f99 - perfectly rational beings (as ASIs are proposed to be)
might be better analysed through above than humans

\subsection{The design space of possible
minds}\label{the-design-space-of-possible-minds}

We are familiar with humans, a few animals, and now are making fairly
limited synthetic intelligences. Yet this probably doesn't exhaust the
possibilities for intelligence. The `shape' of ASI(s) may impact what
agency they have, their goals, and values, and whatever we hope to pass
on to them needs to be thought about in context of what they `are'.

This is the topic for a paper to be presented in April at
\href{https://www.pomconference.org/pom-aachen-2024/}{POM Aachen}

See comprehensive notes at {[}{[}space\_possible\_minds{]}{]}

\begin{itemize}
\tightlist
\item
  Framework

  \begin{itemize}
  \tightlist
  \item
    agency
  \item
    morphology
  \item
    optimisation's role in finding shapes of minds
  \end{itemize}
\item
  singletons vs groups

  \begin{itemize}
  \tightlist
  \item
    does human diversity tell us something about effecting actions in
    the physical world and remaining flexible to respond to local
    stimuli
  \item
    how much larger or smaller can beings be?
  \end{itemize}
\item
  role of evolution in designing us
\item
  prelim list of sources

  \begin{itemize}
  \tightlist
  \item
    sandberg/eckersley =\textgreater{} see
    {[}{[}space\_possible\_minds\#\^{}de0f4b{]}{]}
  \item
    sloman paper =\textgreater{} see
    {[}{[}space\_possible\_minds\#\^{}d36aeb{]}{]}
  \item
    lem of the Summa Technologiae
  \item
    katherine hayles on non-conscious cognitive assemblages
  \item
    shanahan =\textgreater{} see
    {[}{[}space\_possible\_minds\#\^{}e0a577{]}{]}
  \item
    yampolskiy =\textgreater{} see
    {[}{[}space\_possible\_minds\#\^{}938165{]}{]}
  \end{itemize}
\end{itemize}

\subsection{The view from nowhere}\label{the-view-from-nowhere}

\^{}f89d2e

This is one of the least clear and quagmire-prone elements of the
project - most western philosophy explicitly assumes a human subject -
but this project is predicated on the possibility of a world with no
humans (but with intelligence that we can reasonably foresee i.e.~as
opposed to extraterrestrial aliens which we definitely haven't found) -
how do we assess value in such a future world?

This is important because decisions we make w.r.t. AI/ASI may need to be
justified from a perspective broader than `what is good for a particular
group of humans' or `what is good for humanity'. - imagine the Bach
example: - we stop building AGI indefinitely, and as a result are
ill-prepared to deal with various foreseeable risks: another pandemic,
population decline, catastrophic war, and basically fail to reach
technological maturity - it could be argued will have wasted part of the
cosmic endowment (of low entropy energy) - there are lots of assumptions
above (e.g.~is there any connection between AGI and avoiding other
existential risks) - but, if those assumptions are semi-plausible, we
have, in some larger-than-humanity sense, done something irretrievably
`bad' (there may be more concrete utilitarian/longtermist ways of
defining this) - a more straightforward case is that of `AI shutdown':
if we are about to shut down (i.e.~kill) a (sufficiently powerful,
sentient, morally-valuable) AI, we need to have some justification for
this, that somehow incorporates the well-being of that AI - see
\href{https://nickbostrom.com/propositions.pdf}{Bostrom/Shulman} which
discusses possibility that we might cross the threshold set out by where
our created intelligences have moral worth

This feels quite close to moral realism discourse
{[}{[}phd\_stuff\#\^{}8fd708{]}{]}. It also sounds like things Singer,
Wittgenstein, Nagel have argued, but is probably a very different
question.

There might be interesting overlaps with indigenous or non-Western
thought (Buddhism/Shinto) or panpsychism, but I'm not anxious to go too
deep there (in writing anyway).

See {[}{[}view\_from\_nowhere{]}{]} for more particularly Bostrom's idea
of a `cosmic host' sitting atop a hierarchy of morality

\subsection{Propositions upon final
goods}\label{propositions-upon-final-goods}

This section tries to write down a few (admittedly subjective/partial)
possibilities for things that I think humans do that are valuable in a
(view-from-nowhere) sense, i.e.~the universe would be poorer if these
didn't exist. - put another way, what is wrong with a
\href{https://www.lesswrong.com/tag/squiggle-maximizer-formerly-paperclip-maximizer}{`paperclip/squiggle
maximiser'} (or maximiser for varieties of things like utilitronium,
hedonium, etc.)? - I propose art (subject to an appropriate definition)
as one such valuable thing we would like the future to have - This
intersects with \textbf{practice} in development of my project to get
AIs to talk credibly about artwork, discussed
\href{https://tripleampersand.org/work-art-age-cybernetic-criticism/}{here}
in theory and as an \href{http://ukc10014.org/fixedp.html}{artwork} - I
propose the diversity that evolution creates in ecological systems as
another such value - these might have in common `\emph{complexity as a
way of resisting entropy}' (an idea from Bach) which might join up with
{[}{[}\#\^{}f89d2e{]}{]} as an objective value

\subsection{An aside on partiality and
posthumanism}\label{an-aside-on-partiality-and-posthumanism}

\^{}b561c4 What is \textbf{out-of-scope of the project}: - this section
is the only part that talks explicitly about \emph{extinction} and is an
expanded version of this
\href{https://www.lesswrong.com/posts/HaGTQcxqjHPyR9Ju6/unpicking-extinction}{post}
Human extinction has been much discussed and is of limited interest
(mostly through the \emph{option value} argument which might get
developed here) - there are sociological issues surrounding discourse on
post/trans-humanism (e.g.~allegations of eugenicism, universalism,
technophilia, etc. coming from people like Braidotti, Gebru, Torres,
Bender, McQuillan). These are valid criticisms but not the principal
focus of my project (but are acknowledged and summarised) - there are
obviously concerns about race/gender bias and job loss owing to
automation, supposed negative impact on human creativity, etc. These are
not in-scope. - the curious re-animation of accelerationism (via e/acc)
has probably just muddied the waters, and it isn't very clear whether
any lasting insight will come out of it - for now it is out-of-scope
\#\# Simulation arguments

My intuition is that simulation is a rich thread that holds this project
together, both in theory and practice - the genesis of the project, this
post by AI researcher
\href{https://ai-alignment.com/sympathizing-with-ai-e11a4bf5ef6e}{Paul
Christiano}, is depends heavily upon simulation (as a way to `test'
candidate successor civilisations for their `niceness') - thinking about
simulated worlds has weird implications for {[}{[}\#\^{}2a1f99{]}{]} and
{[}{[}\#\^{}f89d2e{]}{]} - the possibility that we will soon be
simulating agentic intelligences might be anthropic evidence regarding
our own `base' reality (per
\href{https://simulation-argument.com/}{Bostrom}) - as a practical
matter, much AI research uses simulations and
\href{http://ukc10014.org/exorbp.html}{my}
\href{http://ukc10014.org/sicul.html}{artistic}
\href{http://ukc10014.org/islands2.html}{practice} and
\href{http://ukc10014.org/writing.html}{writing} have long been built
around this topic

\subsection{On practice}\label{on-practice}

blah blah

\end{document}
